% !TeX root = ../papers/nc-p3.tex
\section{二対に分かれる対称境界の排除}
\label{sec:two-pair-boundary}
二つの互いに素な位置対を可逆な正証明書で同一視し、テンソルを
$\operatorname{Sym}^2 V\otimes\operatorname{Sym}^2 W$の元として書く。
この表示は記述用であり、元の窓の巡回順序は保つ。
対をまたぐ特異な正証明書$H:W\to V$があるとし、$A=\operatorname{im}H$、
$D=\ker H$と置く。このとき
\[
 w\in\operatorname{Sym}^2 A\otimes\operatorname{Sym}^2 W+
        \operatorname{Sym}^2 V\otimes\operatorname{Sym}^2 D.                 \tag{*}
\]
実際、一つの$W$因子に$H$を作用させたテンソルは三つの$V$位置に対称であり、
その三位置の全てが$A$に支持される。従って片方の$W$文字を
$D^\perp=\operatorname{im}H^*$の元で縮約すると両$V$文字が$A$に入る。
もう一つの$W$位置についても対称性を使えば$(*)$を得る。
一つずつ異なる種類の因子へ縮約した像を$F\subset V\otimes W$と書く。
対ごとの対称性から四つの混合対の像は順序を揃えれば同じ$F$であり、
$(*)$から$F\subset A\otimes W+V\otimes D$である。

\begin{proposition}[交互に並ぶ二対]\label{prop:alternating-position-pairs}
窓の巡回順序が$V,W,V,W$なら、$0<r=\rk H<4$であるどの交差証明書も
標準AS性を排除する。
\end{proposition}
\begin{proof}
全ての隣接関係空間は$F$またはその反転である。一次空間を位置ごとに
$V/A,W/D$へ落とすと、全関係が消える。従って長さ$4m$の商の次元は
$\bigl(r(4-r)\bigr)^{2m}$となる。
\end{proof}

\begin{proposition}[隣接する二対の四次障害]\label{prop:adjacent-position-pairs}
窓の巡回順序が$V,V,W,W$で、$w$が標準有限Frobenius候補、
$H$が階数2の交差証明書であるとする。このとき、連続純粋関係があるか、
$\dim C_{0,4}\ge41$である。従ってどちらの場合も標準AS性を満たさない。
\end{proposition}
\begin{proof}
$\dim A=\dim D=2$である。対を分ける切断の階数は6であり、
$(*)$からその上界は$\dim\operatorname{Sym}^2 A+\dim\operatorname{Sym}^2 D=6$となる。
等号より$R_0\supset\operatorname{Sym}^2 A$、$R_2\supset\operatorname{Sym}^2 D$である。
実際、$\operatorname{Sym}^2 A$で像側を割った行列は$(\operatorname{Sym}^2 D)^*$からの階数3写像となり、
元の像が6次元なので核に当たる3次元も全て像に入る。転置にも同じ議論を適用する。
$R_0\subset\operatorname{Sym}^2 V$, $R_2\subset\operatorname{Sym}^2 W$でもある。

以後、連続純粋関係がないとする。
$F=R_1$のうち$F_A=F\cap(A\otimes W)$または$F_D=F\cap(V\otimes D)$に入る
非零純テンソルは存在しない。例えば$a\otimes w'\in F_A$なら$a\otimes a\in R_0$と連続し、
$F_D$でも後続の平方と連続する。$2\times4$行列空間の射影階数1部分は
$\mathbb P^1\times\mathbb P^3\subset\mathbb P^7$の4次元部分なので、
射影次元定理から$\dim F_A,\dim F_D\le3$である。

一つの$V$位置を$A^\perp$、一つの$W$位置を任意の汎関数で縮約する像を$J_D$とする。
$(*)$から$J_D\subset F_D$である。一方、その次元は混合切断の転置を使えば
$F\to(V/A)\otimes W$の像次元$6-\dim F_A$である。よって
$\dim F_A=\dim F_D=3$、$J_D=F_D$である。
$J_D$の$V$側の支持は$V$全体である。もし三次元以下なら、
$3\times2$行列の三次元部分空間には射影次元定理で純テンソルが入り、矛盾する。

ここで$(F_A\otimes W)\cap(A\otimes R_2)=0$を示す。
交叉の元は$W_{1,3}$に入り、有限Frobenius性から、先頭$V$である位置0に
ある$\phi\in V^*$を作用させた$w(\phi)$と書ける。
残る$V$が$A$に支持されるので、そこを$A^\perp$で、さらに$W$の一位置を任意に縮約すると零となる。
従って$\phi$は$J_D$の$V$支持を消し、$\phi=0$である。
ゆえに$A\otimes W\otimes W$から$C_{1,4}$への像次元は
$32-3\cdot4-2\cdot6=8$以下となる。
四次Koszul複体の$d_2$を$\operatorname{Sym}^2 A\otimes C_{2,4}$に制限すると、
その定義域は30次元で、像は$A$と今の8次元以下の像のテンソル積に入る。
従ってこの部分から少なくとも14次元のサイクルが生じる。

このうち境界は8次元以下であることを確認する。
$R'_0=R_0/\operatorname{Sym}^2 A$とし、$d_3:W_{0,3}\otimes W\to R_0\otimes C_{2,4}$を
$R'_0\otimes C_{2,4}$へ射影する。$(*)$の最初の項を除いたテンソルは、
先の階数3から$(\operatorname{Sym}^2 D)^*\simeq R'_0$という同型を与える。
従って$W_{0,3}\simeq W^*$を使うと、射影した$d_3$は
$W^*\twoheadrightarrow D^*$に続いて
\[
 D^*\otimes W\longrightarrow
  (\operatorname{Sym}^2 D)^*\otimes ((D\otimes W)/\operatorname{Sym}^2 D)
  \ \subset R'_0\otimes C_{2,4}
\]
を作用させる写像である。ここで$R_2\cap(D\otimes W)=\operatorname{Sym}^2D$を使った。
写像は二次テンソル$z^2$で評価すると$\lambda\otimes y\mapsto\lambda(z)[z\otimes y]$で与えられる。
$W/D$方向では単射で階数4、$y\in D$の部分では$D^*\otimes D$の元を$T$として
$z\mapsto z\wedge Tz$となり、核はスカラーのみで階数3である。
二つの像は補空間の選択で直和となり、射影の階数は7である。
命題\ref{prop:quartic-koszul-obstruction}より$d_3$自体の階数は15なので、
$\operatorname{Sym}^2 A\otimes C_{2,4}$内に入る境界の次元は$15-7=8$である。
従って$\dim H_2\ge14-8=6$、$\dim C_{0,4}\ge35+6=41$となる。
\end{proof}

\begin{theorem}[対称閉包のAS残余では可逆対は高々一つ]\label{thm:no-two-pair-AS-boundary}
$Z_{\rm sym}\setminus S$内の標準AS候補では、六つの無順序位置対のうち
可逆な正証明書を持つものは高々一対である。
\end{theorem}
\begin{proof}
位置を共有する二対なら三因子を同時対称化でき、基点排除によりASの場合は$S$に入る。
共有しない二対があるなら、それらをまたぐ可逆証明書も同じ理由で存在しない。
系\ref{cor:two-pairs-transfer-ranks}より両方向で階数2の交差証明書を取れる。
二対の巡回順序は上の二命題で尽くされ、定理\ref{thm:no-consecutive-pure-relations}と合わせて
いずれもAS性に反する。よって主張を得る。
\end{proof}
$Z_{\rm sym}$の外では、向かい合う二対に正証明書を持つ符号付きOreの既知AS例と両立する。

\subsection{二対二の三つの境界因子の実現}
ここでは排除条件を満たす境界が実在することと、その余次元を確認する。
$V=A\oplus B$, $W=C\oplus D$と二次元ずつに分け、
$H:C\simeq A$, $K:B\simeq D$を用いて$A,C$を色$r$、
$B,D$を色$s$の共通二次元空間として扱う。
二対内の対称性とこの二つの交差証明書の条件を満たす模型を次で作る。
$V$側に$s$文字があり同時に$W$側に$r$文字がある係数を零とし、
他の係数には同じ色の文字同士で対称な次の二重次数の形式を入れる。
\[
 \alpha_{4,0},\quad\beta_{3,1},\quad\gamma_{2,2},\quad
 \delta_{1,3},\quad\eta_{0,4}.
\]
添字は二色の個数である。係数は対称テンソルの文字座標で取り、多項式の多項係数を掛けない。
$rr,rs,ss$の対の空間の次元を重複を除いて$3,4,3$とする。
$\alpha,\eta$の二対二行列を$A_0,E_0$、$\beta,\delta$の$rr|rs,rs|ss$行列を$B_0,D_0$とする。
$\gamma$には二つの読み方$G_\gamma$ ($rr|ss$, $3\times3$)、
$C_\gamma$ ($rs|rs$, $4\times4$)がある。後者は対称行列全体の中の超平面
\[
 \sigma(C)=C_{00,11}-C_{01,10}=0
\]
の任意の元を一意に表す。混合切断では$sr$の行・列が零で、$rr,ss$内の反対称成分も零である。
重複を除くと混合切断は
\[
 \begin{pmatrix}A_0&B_0&0\\ B_0^{\mathsf T}&C_\gamma&D_0\\
                 0&D_0^{\mathsf T}&E_0\end{pmatrix}.
\]
$A_0,E_0$可逆のとき、階数6となる条件は
$C_\gamma=B_0^{\mathsf T}A_0^{-1}B_0+D_0E_0^{-1}D_0^{\mathsf T}$である。
右辺を$T$と書く。$\sigma(T)=0$が$\alpha,\beta,\delta,\eta$に対する一条件であり、
成立すれば$\gamma$は一意である。

\begin{theorem}[二対二配置による三つの追加境界因子]\label{thm:three-two-pair-divisors}
四位置の二対への三つの分割ごとに、$Z_{\rm sym}\setminus S$には
その二対だけが一般点で可逆正証明書を持つ68次元既約成分が存在する。
これらと定理\ref{thm:four-boundary-divisors}の四成分は全て異なる。
各成分の一般点は有限Frobenius非Koszulであり、$\mathfrak X_6$のある開近傍で
標準AS点は全て$S$に属する。
\end{theorem}
\begin{proof}
まず具体的な非空開チャートを定める。二進添字の和の昇順で
\[
\begin{aligned}
 \alpha&=(1,0,2,1,3), &\eta&=(2,1,3,0,5),\\
 \beta&=((2,3,5,7),(11,13,17,19)), &\delta&=((2,1,0,3),\,\lambda(2,4,1,-1)).
\end{aligned}
\]
を対照点に使う。$\delta$の第二列ベクトルのスカラー$\lambda$に関し$\sigma(T)$は一次である。
実際、$\sigma$の二つの成分はいずれも異なる$r$文字から一つずつ取る。
対照点での傾きは$-54$、根は$\lambda=-23/162$となる。
一般にもこの第二ベクトルの一つの非零座標をスカラーとして分離し、
その方向を射影チャートで変化させれば、傾きが非零な開部分で一意に解ける。
よって$5+8+8+5-1=25$次元の既約な係数チャートを得る。
$\gamma$を加えた五形式全体に一つのスカラーを掛ける自由を除けば24次元である。

このチャートの有限Frobenius部分が$Z_{\rm sym}$に入ることを示す。
完全対称四テンソルの二重次数成分を順に
$\alpha,\beta,\gamma(t),t^{-1}\delta(t),t^{-2}\eta$とする。
その中間縮約の$ss$の行・列に$t$を掛けると、上のブロック表示の二つの零を
$tG_{\gamma(t)}$とその転置で置き換えたものになる。
階数6を保つSchur補の方程式は
\[
 C_{\gamma(t)}=
 (B_0^{\mathsf T}\ \ D_0(t))
 \begin{pmatrix}A_0&tG_{\gamma(t)}\\tG_{\gamma(t)}^{\mathsf T}&E_0\end{pmatrix}^{-1}
 \binom{B_0}{D_0(t)^{\mathsf T}} .
\]
$t=0$での未知数を$\gamma$の9座標と上の$\lambda$とする。
対称$4\times4$行列への線型化は、$\gamma$から$\ker\sigma$へ同型であり、
残る一方向は選んだ非零の傾きとなる。従ってJacobianは可逆で、
次数ごとにこの十元連立一次方程式を解いて形式的な持ち上げを得る。
最初の二位置の$s$文字を$t$倍すれば全係数は$k[[t]]$に入り、特殊化が指定した模型になる。
実際、$s$文字をできるだけ後の二位置に置く成分だけが最小付値を持つ。
一般ファイバーは完全対称テンソルの四因子基底変更で、三分割の階数は6以下である。
有限Frobenius条件はこの三階数切断内で開なので、特殊点がそれを満たせば一般ファイバーも満たす。
完全対称な有限Frobeniusテンソルの認識定理より一般ファイバーは$S$に属する。
従って所与の特殊点は$Z_{\rm sym}$に入る。

この対照点では、一因子階数は全て4、三切断の階数は全て6である。
隣接二対の順序でも、一対を交互に配置する順序でも、四始点全てで
$(\dim W_{i,3},\dim W_{i,4})=(4,1)$である。
全12個の向き付き正symmetrizer空間は一次元であり、対の内部で階数4、対をまたぐと階数2である。
従ってこれらの次元1条件は上の族の非空開部分で保たれる。
三切断の射影接次元の有理数検算は69である。
数値による近似は使わず、\path{computations/verify_two_pair_boundary.py}で関係挿入の核、
行列の階数とSchur方程式の傾きを計算する。

位置のラベルを固定して表示の自由度を数える。二対の各内部で一方の因子の基底を合わせる
射影変換が二つで30次元。共通空間$V,W$内の二平面$A,D$の選択が$4+4$次元、
射影同型$W/D\simeq A$, $V/A\simeq D$の選択が$3+3$次元である。
この14次元の基底データのGrassmann・行列チャート内で補空間と基底を正則に選べる。
係数を合わせて$30+14+24=68$次元の既約な表示族となる。
一般点から、二対内の一次元可逆証明書で二対をそれぞれ正規化し、
両向きの交差証明書の射影類を回収すれば$A,D$と二つの射影同型が決まる。
固定した局所基底では五形式も共通スカラー以外一意である。よって一般ファイバーは零次元である。
この族は$S$と交わらない一般部分を持つので、その68次元閉包は69次元成分の境界因子である。

四位置を置換すると他の二分割の因子も得られる。巡回順序を変えた場合の有限Frobenius性は
上の二順序の検算と開性により確かめられる。一般点で可逆な位置対が指定した二対だけであることから、
それら同士および三因子対称の四つの境界因子との区別がつく。
対照点で$\mathfrak X_6$は69次元成分を含み接次元も69なので滑らかであり、近傍を
$Z_{\rm sym}$のみに取れる。その中で二対の可逆証明書は、次元1と可逆性の開性で残る。
定理\ref{thm:no-two-pair-AS-boundary}より、この近傍にある標準AS点は$S$に属する。
各因子の稠密開部分について主張が従う。七つで全境界を尽くすかは未決定である。
\end{proof}

対照点の隣接配置では$F_A,F_D$の二行または二列の行列束を調べると、
どちらも四つの最大minorの最大公約式が1であり、純テンソルを含まない。
$F\subset A\otimes W+V\otimes D$内の純テンソルはどちらかに入るため、混合関係には純テンソルがない。
このとき四次経路次元は$(41,42,41,42)$で、隣接配置の下界41は始点0で達成される。
同じ係数を交互に配置すると四次経路次元は全始点で35だが、自由な交互商による指数成長がある。

\begin{corollary}[二対二境界の局所方程式]\label{cor:two-pair-transverse-coordinate}
上の各対照点の近くで、対をまたぐ両向きの正証明書直線の局所生成元を$H,K$と取ると、
$HK=KH=zI$のスカラー$z$が滑らかな被約境界の局所方程式である。
両行列式はともにこの因子に沿って二重に消える。
\end{corollary}
\begin{proof}
二対の内部を同一視した局所座標を使う。
$Z_{\rm sym}$の稠密開部分$S$の上で$HK,KH$は同じスカラー行列であり、
この等式は正則な局所生成元に対して近傍全体に延びる。
対内の可逆性も保たれるため、$z\ne0$は四因子全てを対称化できることと同値である。
前定理の形式曲線では最初の二位置の色$s$を$t$倍するので、
$H,K$は単元を除いて$\operatorname{diag}(1,1,t,t)$と
$\operatorname{diag}(t,t,1,1)$である。従って$z$は曲線上で単元倍の$t$となり、
滑らかな周囲で$dz\ne0$を得る。$z=0$で両行列は階数2
（階数2以上は開であり、積が零）なので、単元二次minorによる掃き出しで
両行列式が$z^2$で割り切れる。積は$z^4$であり、両方の消滅次数は正確に2である。
\end{proof}

\begin{corollary}[二対を持つAS点の交差証明書空間]\label{cor:two-pair-cross-space-zero}
$Z_{\rm sym}$への所属を仮定せず、標準AS対象が互いに素な二つの位置対で
可逆な正証明書を持つとする。$S$に属さなければ、異なる対の位置間の
正symmetrizer空間は両向きとも零である。
\end{corollary}
\begin{proof}
非零の交差証明書$H$が可逆なら、二対の同一視と合わせて四因子を対称化できる。
階数1は系\ref{cor:no-rank-one-into-positive-pair}で排除され、
階数2はいずれの巡回配置でも上の二命題で排除される。
階数3では$H$を$\operatorname{diag}(I_3,0)$に取る。対象二因子について、
残る因子の全縮約行列$M$は$MH^{\mathsf T}$が対称であるという条件を満たす。
従って$M$の左下$1\times3$ブロックは零である。
逆向きに$K=\operatorname{diag}(0,0,0,1)$を取れば、$KM$は右下の一要素しか持たず対称である。
よって$K$は逆向きの階数1の正証明書となり、その標的位置も可逆対に入っているので矛盾する。
\end{proof}
